% Summary of the permeation chapter as a resistance network. Fills encode the tritium
% concentration, which falls along every pathway; R3 is drawn as a variable resistance because
% the sweep gas chemistry is what sets it, and the wall is drawn as a tank because when R3 is
% large it stores instead of conducting.
% Requires in the preamble: \usepackage{tikz}
%   \usetikzlibrary{positioning, fit, backgrounds, arrows.meta}
\begin{tikzpicture}[
    >=Stealth,
    wire/.style={line width=1.1pt, black!65},
    res/.style={draw=black!75, fill=white, line width=0.8pt,
                minimum width=1.5cm, minimum height=0.62cm, inner sep=1pt},
    tank/.style={draw=black!70, line width=0.9pt, rounded corners=2pt},
    sink/.style={draw=black!55, line width=0.7pt, rounded corners=2pt,
                 minimum width=2.1cm, minimum height=0.95cm, align=center, font=\footnotesize},
    note/.style={font=\scriptsize, black!60, align=center},
    val/.style={font=\scriptsize, align=center},
    font=\small,
]
% ---- the two reservoirs -------------------------------------------------------------- %
\filldraw[tank, fill=red!45] (0,0) rectangle (2.5,3.0);
\node[align=center] at (1.25,1.95) {\textbf{salt}};
\node[note, align=center] at (1.25,1.25) {the whole bred\\inventory};

\filldraw[tank, fill=orange!45] (6.0,0) rectangle (7.1,3.0);
\node[rotate=90] at (6.55,1.5) {\textbf{wall}};
\node[note, anchor=south east, align=right] at (7.15,3.12) {$R_2$, diffusion\\in inconel, fast};

% ---- sparging, straight out of the salt ------------------------------------------------ %
\draw[wire, ->] (1.25,3.0) -- (1.25,4.15);
\node[sink, fill=black!4, minimum width=2.7cm] at (1.25,4.68)
    {extracted by sparging, $\tau$};

% ---- R1: salt to wall, the one that decides how much permeates -------------------------- %
\draw[wire] (2.5,1.5) -- (6.0,1.5);
\node[res, fill=red!12] at (4.25,1.5) {$R_1$};
\node[val, anchor=south] at (4.25,1.88) {$1/k_\mathrm{eff}$, $k_\mathrm{eff}=k_m/\mathrm{PRF}$};

% ---- the divider: R3 to the OV gas, R4 to the room -------------------------------------- %
\draw[wire] (7.1,2.35) -- (8.05,2.35);
\draw[wire, ->] (9.55,2.35) -- (10.3,2.35);
\node[res, fill=orange!12] at (8.8,2.35) {$R_3$};
\draw[->, black!75, line width=0.7pt] (8.15,1.98) -- (9.45,2.72);   % variable resistance
\node[val, anchor=south] at (8.8,2.98) {sweep gas chemistry};

\draw[wire] (7.1,0.65) -- (8.05,0.65);
\draw[wire, ->] (9.55,0.65) -- (10.3,0.65);
\node[res, fill=orange!12] at (8.8,0.65) {$R_4$};
\node[val, anchor=north] at (8.8,0.30) {spacer ribs, set by $\chi$};

\node[sink, fill=black!4] at (11.45,2.35) {collected\\in the OV gas};
\node[sink, fill=black!4, draw=red!70!black] at (11.45,0.65)
    {\textcolor{red!70!black}{leaked}\\to the environment};

% ---- the two statements the network makes ----------------------------------------------- %
\node[note, anchor=north, align=center] at (4.25,-0.30)
    {$R_1$ \textbf{dominates}: it alone sets \emph{how much}\\
     permeates, whatever the sweep gas is};
\node[val, anchor=north west, align=left] at (9.55,-0.30)
    {\textbf{share that leaks}\\
     He only \quad $\sim100\ \%$\\
     $+$\,\ce{H2O} \quad $35-45\ \%$\\
     $+3.5\ \%$\,\ce{H2} \quad $0.3\ \%$};
\node[note, anchor=north west, align=left] at (0,-1.55)
    {$R_3$ large (He only): the wall stops passing tritium on and\\
     fills up instead, holding $84\ \%$ of the bred inventory after 4 days};

% ---- the concentration gradient the fills encode ---------------------------------------- %
\shade[left color=red!45, right color=black!4] (0,-2.85) rectangle (2.4,-2.58);
\draw[black!55, line width=0.5pt] (0,-2.85) rectangle (2.4,-2.58);
\node[note, anchor=west] at (2.55,-2.72)
    {fill colour: tritium concentration, and with it the potential to leak};
\end{tikzpicture}
